\documentclass[12pt]{article}
\usepackage[utf8]{inputenc}
\usepackage[russian]{babel}
\usepackage{graphicx}
\usepackage{array}
\usepackage[a4paper, left=1cm, right=1cm, top=1.27cm, bottom=1.27cm]{geometry}
\usepackage{setspace}
\usepackage{makecell}
\usepackage{caption}

\usepackage{minted}
\usepackage{xcolor}
\usepackage{geometry}
\usepackage{amsmath}
\usepackage{amssymb}
\geometry{margin=2cm}


\begin{document}

\begin{flushright}
    Пахомов Владимир Юрьевич, 467036
\end{flushright}

\begin{center}
    \textbf{Домашнее задание №2}
\end{center}

\begin{flushleft}
    \textbf{Задание 1} \\
    Условие: \\
    Даны три части секрета: $(1, 567), (2, 890), (3, 1023)$. Восстановите секрет $S$, если $p = 1237$ \\
    
    \vspace{20pt}
    \textbf{Решение}: \\
    Используем для вычисления S интерполяционный многочлен Лагранжа: \\
    \[ S = \sum_{j=1}^{3} y_j \cdot L_j \; \; (mod \; p) \]
    \[ L_j = \Pi_{i\ne j}\frac{-x_i}{x_j - x_i} \]
    Считаем: \\
    \[ L_1 = \frac{-2}{1 - 2} \cdot \frac{-3}{1 - 3} = 3 \]
    \[ L_2 = \frac{-1}{2 - 1} \cdot \frac{-3}{2 - 3} = -3 \]
    \[ L_3 = \frac{-1}{3 - 1} \cdot \frac{-2}{3 - 2} = 1 \]
    \[ S = 567 \cdot 3 + 890 \cdot (-3) + 1023 \cdot 1 \; \; (mod \; 1237) \]
    \[ S = 54 \]
\end{flushleft}


\newpage


\begin{flushleft}
    \textbf{Задание 2} \\
    Условие: \\
    Докажите, что $(2p-1)! - p$ делится на $p^2$, если $p$ - простое\\
    
    \vspace{20pt}
    \textbf{Решение}: \\
    Нужно доказать, что $(2p-1)! \equiv p \; \; (mod \; p^2) $ \\
    Разложим фактоирал:
    \[ (2p-1)! = 1 \cdot 2 \dots (p-1) \cdot p \cdot (p+1) \dots (2p-1) = p \cdot (p-1!) \cdot (p+1) \cdot (p+2) \dots (2p - 1) \]
    Обозначим все проивзедение после $p$ за $x$, тогда: \\
    \[(2p-1)! = px \]
    Теперь нужно доказать, что $px - p$ делится на $p^2$, но достаточно доказать, что
    \[ x \equiv 1 \; \; (mod \; p) \]
    Заметим, что $p + k \equiv k \; \; (mod \; p), \; \forall k$ \\
    Получается, что 
    \[(p+1) \cdot (p+2) \dots (2p - 1) \equiv (p-1)! \; \; (mod \; p) \]
    Подставляя обратно в $x$:
    \[ x \equiv (p-1)!^2 \; \; (mod \; p) \]
    По теореме Вильсона $(p-1)! \equiv -1 \; \; (mod \; p)$ \\
    Следовательно
    \[ x \equiv (-1)^2 \; \; (mod \; p) \; \Rightarrow \; x \equiv 1 \; \; (mod \; p) \]
    \[ x = kp + 1 \]
    \[(2p-1)! = p \cdot x = p(kp + 1) = kp^2 + p \]
    \[(2p-1)! - p = kp^2 \]
    Так как $k \; -$ целое, то $kp^2$ делится на $p^2$, следовательно и $(2p-1)!-p$ тоже будет делиться на $p^2$ 

\end{flushleft}


\newpage


\begin{flushleft}
    \textbf{Задание 3} \\
    Условие: \\
    Составить таблицы сложения и умножения: \\
    \[ \mathbb{Z}_3[x]/(x^2 + x + 2) \]
    
    \vspace{20pt}
    \textbf{Решение}: \\
    Данное поле состоит из полиномов, степени меньше 2. Коэффициенты $\{0, 1, 2\}$. 
    Каждый полином имеет вид $ax + b$. То есть всего может быть $3 \cdot 3 = 9$  элементов:
    \[\{0, \; 1, \; 2, \; x, \; x+1, \; x+2, \; 2x, \; 2x + 1, \; 2x + 2\} \]
    Операции совершаются по модулю $x^2 + x + 2$. Выразим $x^2$:
    \[ x^2 + x + 2 \; \Rightarrow \; x^2 = -x - 2 \; \Rightarrow \; x^2 = 2x + 1 \; \; (mod \; 3) \]
    Таблица сложения: \\
    Пример: $(2x+1) + (x+2) = 3x + 3 = 0 \; \; (mod \; 3)$ \\ 

    \begin{table}[h]
        \centering
        \renewcommand{\arraystretch}{1.2}
    \begin{tabular}{c|c|c|c|c|c|c|c|c|c}
        + & 0 & 1 & 2 & x & x+1 & x+2 & 2x & 2x+1 & 2x+2 \\
        \hline
        0 & 0 & 1 & 2 & x & x+1 & x+2 & 2x & 2x+1 & 2x+2 \\
        1 & 1 & 2 & 0 & x+1 & x+2 & x & 2x+1 & 2x+2 & 2x \\
        2 & 2 & 0 & 1 & x+2 & x & x+1 & 2x+2 & 2x & 2x+1 \\
        x & x & x+1 & x+2 & 2x & 2x+1 & 2x+2 & 0 & 1 & 2 \\
        x+1 & x+1 & x+2 & x & 2x+1 & 2x+2 & 2x & 1 & 2 & 0 \\
        x+2 & x+2 & x & x+1 & 2x+2 & 2x & 2x+1 & 2 & 0 & 1 \\
        2x & 2x & 2x+1 & 2x+2 & 0 & 1 & 2 & x & x+1 & x+2 \\
        2x+1 & 2x+1 & 2x+2 & 2x & 1 & 2 & 0 & x+1 & x+2 & x \\
        2x+2 & 2x+2 & 2x & 2x+1 & 2 & 0 & 1 & x+2 & x & x+1 \\
    \end{tabular}
    \end{table}

        
    Таблица умножения: \\
    Пример: $(2x+1) \cdot (x+2) = 2x^2 + x + 4x + 2 = 2(2x + 1) + 5x + 2 = 9x + 4 = 1 \; \; (mod \; 3) $ \\

    \begin{table}[h]
        \centering
        \renewcommand{\arraystretch}{1.2}
    \begin{tabular}{c|c|c|c|c|c|c|c|c|c}
    $\times$ & 0 & 1 & 2 & x & x+1 & x+2 & 2x & 2x+1 & 2x+2 \\
        \hline
        0 & 0 & 0 & 0 & 0 & 0 & 0 & 0 & 0 & 0 \\
        1 & 0 & 1 & 2 & x & x+1 & x+2 & 2x & 2x+1 & 2x+2 \\
        2 & 0 & 2 & 1 & 2x & 2x+2 & 2x+1 & x & x+2 & x+1 \\
        x & 0 & x & 2x & 2x+1 & 1 & x+1 & x+2 & 2x+2 & 2 \\
        x+1 & 0 & x+1 & 2x+2 & 1 & x+2 & 2x & 2x+2 & 2 & x \\
        x+2 & 0 & x+2 & 2x+1 & x+1 & 2x & 2x+2 & 2 & 1 & x \\
        2x & 0 & 2x & x & x+2 & 2x+2 & 2 & 2x+1 & x+1 & 1 \\
        2x+1 & 0 & 2x+1 & x+2 & 2x+2 & 2 & 1 & x+1 & 2 & x \\
        2x+2 & 0 & 2x+2 & x+1 & 2 & x & x & 1 & x & x+2 \\
    \end{tabular}
    \end{table}
\end{flushleft}


\newpage


\begin{flushleft}
    \textbf{Задание 4} \\
    Условие: \\
    \[g=5, \; a=43, \; p = 257(=2^16 + 1) \]
    Найти: $ x: \; g^x \equiv a \; \; (mod \; p) $ \\
    
    \vspace{20pt}
    \textbf{Решение}: \\
    \[5^x \equiv 43 \; \; (mod \; 257) \]
    \[ x = q_0 + 2q_1 + 4q_2 + _{\dots \dots} + 128q_7 \]
    \[ a_0^{\frac{p-1}{2}} \; \; (mod \; p) \]

    \vspace{20pt}
    Нахождение $q_0$: \\
    \[ 43^{128} \pmod{257} \]
    $43^2 = 1849 = 7 \cdot 257 + 50 \equiv 50 $\\
    $43^4 \equiv 50^2 = 2500 \equiv -70$ \\
    $43^8 \equiv (-70)^2 = 4900 \equiv 17$ \\
    $43^{16} \equiv 17^2 = 289 \equiv 32$ \\
    $43^{32} \equiv 32^2 = 1024 \equiv -4$ \\
    $43^{64} \equiv (-4)^2 = 16$ \\
    $43^{128} \equiv 16^2 = 256 \equiv -1$ \\
    Так как результат $-1$, то бит \textbf{$q_0 = 1$}.

    \vspace{20pt}
    Нахождение $q_1$: \\
    $a_1 = a_0 \cdot g^{-q_0} = 43 \cdot 5^{-1}$.
    Найдём $5^{-1} \pmod{257}$. Используя расширенный алгоритм Евклида:
    \[ 5 \cdot 103 = 515 = 2 \cdot 257 + 1 \implies 5^{-1} \equiv 103 \]
    \[ a_1 = 43 \cdot 103 = 4429 = 17 \cdot 257 + 60 \equiv 60 \]
    Проверяем $a_1^{\frac{p-1}{4}} = 60^{64} \pmod{257}$.
    Вычислим степени 60:
    \[ 60^2 = 3600 = 14 \cdot 257 + 2 \equiv 2\]
    \[60^{64} = (60^2)^{32} = 2^{32}\]
    Так как $2^{16} \equiv 1$, то $2^{32} \equiv 1$. Результат $1$, значит \textbf{$q_1 = 0$}.

    \vspace{20pt}
    Нахождение $q_2$: \\
    Так как $q_1=0$, значение $a_2 = a_1 = 60$.
    Проверяем $a_2^{\frac{p-1}{8}} = 60^{32} \pmod{257}$.
    \[60^{32} = (60^2)^{16} = 2^{16} \equiv 1 \]
    Результат $1$, значит \textbf{$q_2 = 0$}.

    \vspace{20pt}
    Нахождение $q_3$ \\
    Так как $q_2=0$, значение $a_3 = 60$.
    Проверяем $a_3^{\frac{p-1}{16}} = 60^{16} \pmod{257}$.
    \[60^{16} = (60^2)^{8} = 2^8 = 256 \equiv -1\]
    Результат $-1$, значит \textbf{$q_3 = 1$}.

    \vspace{20pt}
    Нахождение $q_4$: \\
    Нам нужно $a_4 = a_3 \cdot (g^8)^{-1}$.
    Вычислим $g^8 = 5^8$:
    \[5^2 = 25, \quad 5^4 = 625 \equiv 111, \quad 5^8 \equiv 111^2 = 12321 \equiv -15\]
    Обратный к $-15$. Так как $15 \cdot 120 = 1800 = 7 \cdot 257 + 1$, то $15^{-1} = 120$.
    Тогда $(-15)^{-1} \equiv -120 \equiv 137$.
    \[a_4 = 60 \cdot 137 = 8220 = 31 \cdot 257 + 253 \equiv -4\]
    Проверяем $a_4^{\frac{p-1}{32}} = (-4)^8 \pmod{257}$.
    \[ (-4)^8 = 4^8 = (2^2)^8 = 2^{16} \equiv 1 \]
    Результат $1$, значит \textbf{$q_4 = 0$}.

    \vspace{20pt}
    Нахождение $q_5$: \\
    Так как $q_4=0$, $a_5 = a_4 = -4$.
    Проверяем $a_5^{\frac{p-1}{64}} = (-4)^4 \pmod{257}$.
    \[ (-4)^4 = 256 \equiv -1 \]
    Результат $-1$, значит \textbf{$q_5 = 1$}.

    \vspace{20pt}
    Нахождение $q_6$ и $q_7$: \\
    $a_6 = a_5 \cdot (g^{32})^{-1}$.
    Мы знаем $g^8 \equiv -15$.
    \[g^{16} \equiv (-15)^2 = 225 \equiv -32\]
    \[g^{32} \equiv (-32)^2 = 1024 \equiv -4\]
    Обратный к $-4$ равен $64$ (так как $-4 \cdot 64 = -256 \equiv 1$).
    \[a_6 = (-4) \cdot 64 = -256 \equiv 1\]
    Так как текущее значение равно $1$, все оставшиеся биты равны $0$.
    \[q_6 = 0, \quad q_7 = 0\]

    Соберем число $x$ из полученных битов $(q_7 \; q_6 \; q_5 \; q_4 \; q_3 \; q_2 \; q_1 \; q_0)_2 = (00101001)_2$:
    \[ x = 1 \cdot 2^0 + 0 \cdot 2^1 + 0 \cdot 2^2 + 1 \cdot 2^3 + 0 \cdot 2^4 + 1 \cdot 2^5 \]
    \[ x = 1 + 8 + 32 = 41 \]

\end{flushleft}


\newpage


\begin{flushleft}
    \textbf{Задание 5} \\
    Условие: \\
    Найти порядок 2 в $(\mathbb{Z}/29\mathbb{Z})^*$ \\
    
    \vspace{20pt}
    \textbf{Решение}: \\
    Нужно найти мультипликативный порядок элемента $a = 2$ по модулю $n=29$. То есть нужно найти такое минимальное $k$, что: 
    \[ 2^k \equiv 1 \; \; (mod \; 29) \]
    Так как $29 \; -$ простое, то порядок группы равен
    \[ \phi(29) = 29 - 1 = 28 \]
    По теореме Лагранжа порядок числа 2 должен делить 28. Возможные порядки $\{1, 2, 4, 7, 14, 28\}$. Проверяем делители: \\
    $2^1 = 2 \neq 1 $ \\
    $2^2 = 4 \neq 1 $ \\
    $2^4 = 16 \neq 1 $ \\
    $2^7 = 128 \equiv 12 \neq 1 $ \\
    $2^{14} = (2^7)^2 \equiv 12^2 = 144 \equiv 28 \equiv -1 \neq 1 $ \\
    $2^{28} =(2^{14})^2 \equiv (-1)^2 = 1 $ \\
    Следовательно порядок 2 в $(\mathbb{Z}/29\mathbb{Z})^*$ равен 28\\

\end{flushleft}

\end{document}

